%
% pcm-stream.tex
%
% Z specification of the streaming PCM sink (src/punt_vox/voxd/pcm_stream.py,
% PL-1298): the phases of one streamed utterance and the byte accounting the
% tests in tests/test_pcm_stream.py assert by name.
%

\documentclass[a4paper,10pt,fleqn]{article}
\usepackage[margin=1in]{geometry}
\usepackage{fuzz}

\begin{document}

\title{The Streaming PCM Sink: A Z Specification}
\author{Punt Labs, for PL-1298}
\date{November 2026}
\maketitle

\section{Introduction}

A \texttt{PcmStreamSink} carries one agent utterance from the network to the
output device. The event loop feeds chunks, finishes the response, flushes on
an ElevenLabs \texttt{interruption}, and closes the sink; the device thread
pulls 20~ms blocks. Pending audio lives in one locked buffer, so each feed,
pull and flush is atomic and the operations below interleave in any order.

The model counts bytes. It names eight properties, and the tests carry the
same names:

\begin{enumerate}
  \item[P1] \emph{Conservation}: every byte fed is played, dropped or still
    pending, exactly once.
  \item[P2] \emph{Flush is exact}: once flushed, nothing is pending and the
    device pulls no more of the utterance.
  \item[P3] \emph{Closed streams never accumulate audio}: no operation adds to
    a closed sink.
  \item[P4] \emph{Stop is idempotent}: a flush after a flush, a close or a
    play-out changes nothing.
  \item[P5] \emph{Play-out drains}: close lets the device play its last block
    only when the response finished and nothing is pending; otherwise it
    aborts.
  \item[P6] \emph{A dry spell opens only on a streaming sink that already
    played audio}: silence before the first audio, or after the response
    finished or was flushed, is not a dry spell.
  \item[P7] \emph{Shaping moves no byte}: the fades change sample values,
    never \emph{fed}, \emph{played}, \emph{dropped} or \emph{pending}; every
    operation on a spell is $\Xi$ on \emph{Sink}.
  \item[P8] \emph{Starvation is lateness, not length}: a spell that ends
    with audio of the same response is starvation at any length; a spell that
    ends at a response boundary, or with the response (finish, flush) or the
    stream, is idle and counted apart.
\end{enumerate}

\section{Types}

\begin{zed}
  PHASE ::= streaming | finished | interrupted | closed \\
  ZBOOL ::= ztrue | zfalse
\end{zed}

\section{State}

\begin{schema}{Sink}
  phase : PHASE \\
  fed, played, dropped, pending : \nat \\
  drained : ZBOOL
\where
  fed = played + dropped + pending \\
  phase \in \{ interrupted, closed \} \implies pending = 0 \\
  drained = ztrue \implies pending = 0
\end{schema}

The first predicate is P1. The second is the state half of P2 and P3.

\begin{schema}{InitSink}
  Sink'
\where
  phase' = streaming \\
  fed' = 0 \\
  played' = 0 \\
  dropped' = 0 \\
  pending' = 0 \\
  drained' = zfalse
\end{schema}

\section{Operations on the event loop}

\begin{schema}{Feed}
  \Delta Sink \\
  n? : \nat_1
\where
  phase = streaming \\
  fed' = fed + n? \\
  pending' = pending + n? \\
  phase' = phase \\
  played' = played \\
  dropped' = dropped \\
  drained' = drained
\end{schema}

Audio still in flight after an interruption is counted and dropped (P1, P2).

\begin{schema}{FeedAfterFlush}
  \Delta Sink \\
  n? : \nat_1
\where
  phase = interrupted \\
  fed' = fed + n? \\
  dropped' = dropped + n? \\
  pending' = pending \\
  phase' = phase \\
  played' = played \\
  drained' = drained
\end{schema}

Feeding a finished or closed sink has no operation: the implementation raises
\texttt{ValueError} (P3).

\texttt{Finish} drops a trailing half sample, which a network chunk boundary
can leave, so the stream can still play out.

\begin{schema}{Finish}
  \Delta Sink \\
  trim? : \nat
\where
  phase = streaming \\
  trim? \leq 1 \\
  trim? \leq pending \\
  pending' = pending - trim? \\
  dropped' = dropped + trim? \\
  phase' = finished \\
  fed' = fed \\
  played' = played \\
  pending' = 0 \implies drained' = ztrue \\
  pending' \neq 0 \implies drained' = drained
\end{schema}

A response that ends with nothing left to play is played out at once.

\begin{schema}{Flush}
  \Delta Sink
\where
  phase \in \{ streaming, finished \} \\
  drained = zfalse \\
  dropped' = dropped + pending \\
  pending' = 0 \\
  phase' = interrupted \\
  fed' = fed \\
  played' = played \\
  drained' = drained
\end{schema}

\begin{schema}{FlushNoOp}
  \Xi Sink
\where
  phase \in \{ interrupted, closed \} \lor drained = ztrue
\end{schema}

\texttt{Flush} and \texttt{FlushNoOp} together are total: P4.

\begin{schema}{Close}
  \Delta Sink \\
  drain! : ZBOOL
\where
  phase \neq closed \\
  (phase = finished \land pending = 0) \implies drain! = ztrue \\
  \lnot (phase = finished \land pending = 0) \implies drain! = zfalse \\
  dropped' = dropped + pending \\
  pending' = 0 \\
  phase' = closed \\
  fed' = fed \\
  played' = played \\
  drained' = drained
\end{schema}

\texttt{drain!} is P5: \texttt{ztrue} stops the device after its last block,
\texttt{zfalse} aborts it.

\section{Operations on the device thread}

\begin{schema}{Pull}
  \Delta Sink \\
  k? : \nat_1
\where
  phase \in \{ streaming, finished \} \\
  k? \leq pending \\
  played' = played + k? \\
  pending' = pending - k? \\
  phase' = phase \\
  fed' = fed \\
  dropped' = dropped \\
  drained' = drained
\end{schema}

\texttt{Pull} needs \texttt{pending} $> 0$, which an interrupted or closed sink
never has: the operational half of P2.

\begin{schema}{Drain}
  \Delta Sink
\where
  phase \in \{ finished, interrupted \} \\
  pending = 0 \\
  drained = zfalse \\
  drained' = ztrue \\
  phase' = phase \\
  fed' = fed \\
  played' = played \\
  dropped' = dropped \\
  pending' = pending
\end{schema}

\texttt{Drain} is the callback finding the buffer empty after
\texttt{Finish}; it releases waiters. The callback reads the phase without
the loop's cooperation, so a \texttt{Flush} may land between that read and
the drain: \texttt{Drain} is therefore allowed from \texttt{interrupted}
too. Nothing is pending in either case, so the two orders end in the same
counts; only the report's \texttt{interrupted} flag differs.

\section{Dry spells}

A dry spell is the buffer running empty while audio was playing and the
response was still arriving. The device thread fades the last sample down to
zero (a 3~ms tail), plays silence, and fades the first samples of resumed
audio up from where the tail had got to. The same tail also follows a flush
or the end of the response. The fades add no phase to the sink: they are the
sample-level shape of the steps below. What the model adds is the spell's own
state, so that the counts and the length the report carries are defined.

The call's stream stays open between the agent's turns, so the sink cannot
tell a pause from a stall by itself. The caller marks where each response
begins; a spell that ends at such a mark is a turn boundary. In the model the
mark is the input \emph{boundary?}: it is true when the audio that ends the
spell is the first of a new response.

\begin{schema}{Spells}
  open : ZBOOL \\
  missing, starved, idle, silence : \nat
\where
  open = zfalse \implies missing = 0
\end{schema}

\emph{silence} is the total length of the starvations; \emph{missing} is the
length of the open spell.

\begin{schema}{InitSpells}
  Spells'
\where
  open' = zfalse \\
  missing' = 0 \\
  starved' = 0 \\
  idle' = 0 \\
  silence' = 0
\end{schema}

\begin{schema}{Observed}
  Sink \\
  Spells
\where
  open = ztrue \implies played > 0
\end{schema}

The predicate of \texttt{Observed} is P6: a spell is open only on a sink
that has played audio.

A callback that comes up short extends the open spell, or opens one. The
code pulls whole frames, so a half sample may stay pending: \emph{pending}
is at most one byte.

\begin{schema}{Starve}
  \Xi Sink \\
  \Delta Spells \\
  k? : \nat_1
\where
  phase = streaming \\
  played > 0 \\
  pending \leq 1 \\
  open' = ztrue \\
  missing' = missing + k? \\
  starved' = starved \\
  idle' = idle \\
  silence' = silence
\end{schema}

\texttt{Starve} is P6 (the preconditions) and P7 (the $\Xi$~\emph{Sink}).

\begin{schema}{Resume}
  \Xi Sink \\
  \Delta Spells \\
  boundary? : ZBOOL
\where
  open = ztrue \\
  open' = zfalse \\
  missing' = 0 \\
  boundary? = ztrue \implies
    (idle' = idle + 1 \land starved' = starved \land silence' = silence) \\
  boundary? = zfalse \implies
    (starved' = starved + 1 \land silence' = silence + missing \land idle' = idle)
\end{schema}

\texttt{Resume} is audio returning. The length of the spell does not appear
in the choice: P8.

\begin{schema}{EndWithResponse}
  \Xi Sink \\
  \Delta Spells
\where
  open = ztrue \\
  phase \in \{ finished, interrupted, closed \} \\
  open' = zfalse \\
  missing' = 0 \\
  idle' = idle + 1 \\
  starved' = starved \\
  silence' = silence
\end{schema}

\texttt{EndWithResponse} closes a spell the response outlived, on the next
callback after \texttt{Finish} or \texttt{Flush}, or at \texttt{Close}. Nothing
was late, so it is idle (P8).

One callback can end a spell and open the next: the resumed audio runs out
inside the same block.

\begin{zed}
  ResumeThenStarve \defs Resume \semi Starve
\end{zed}

\end{document}
